\documentclass[12pt, a4paper]{article}
\usepackage[utf8]{inputenc}
\usepackage[T1]{fontenc}
\usepackage{amsmath, amssymb}
\usepackage{graphicx}
\usepackage{hyperref}
\usepackage{geometry}
\usepackage{biblatex}
\usepackage{booktabs}
\usepackage{listings}
\usepackage{xcolor}

\geometry{margin=2.5cm}
\addbibresource{bibliography/references.bib}

\title{\textbf{Glass Box Architectures: The Ontogeny of Symbolic Reasoning in Hybrid Neural Networks}\\
\large A Mechanistic Analysis of Transformers, State Space Models, and Mixture of Experts}
\author{Manu (PhD Candidate) \\ Supervisor: Antigravity AI}
\date{2024-2025}

\begin{document}

\maketitle

\begin{abstract}
\noindent Large Language Models (LLMs) are typically deployed as opaque "Black Boxes," where the emergence of higher-order reasoning remains a sub-symbolic mystery. This thesis challenges that paradigm by proposing a "Glass Box" framework, utilizing a Hybrid Architecture (\textit{Cortex-13}) composed of Transformers, Mamba State Space Models, and Mixture of Experts (MoE). By instrumenting this architecture with Sparse Autoencoders (SAEs) and subjecting it to a biologically inspired "Developmental Curriculum" (Innate $\rightarrow$ Social $\rightarrow$ Specialized), we demonstrate that symbolic logic is not an inherent property of the architecture but an emergent phenomenon of information compression. We provide quantitative evidence that reinforcement learning induces a phase transition in feature sparsity ($p < 0.05$), effectively crystallizing continuous activation vectors into discrete boolean logic gates.
\end{abstract}

\tableofcontents
\newpage

\section{1. Introduction}
The field of Artificial Intelligence has largely bifurcated into two camps: the high-performing but opaque Connectionist models (Deep Learning) and the interpretable but brittle Symbolic reasoning systems. The "Neuro-Symbolic" gap remains the central challenge of the decade.

\subsection{1.1 The Interpretability Crisis}
While models like GPT-4 exhibit behavior that mimics reasoning, Mechanistic Interpretability research \cite{GemmaScope2024} suggests they largely rely on polysemantic neurons—units that activate for unrelated concepts (e.g., a neuron firing for both "The color Blue" and "Historical Dates"). This superposition makes "Glass Box" analysis nearly impossible with standard tools.

\subsection{1.2 The Ontogenetic Hypothesis}
We propose that the "Black Box" nature is a transient state of immature models. Drawing from developmental cognitive neuroscience, we hypothesize that symbolic clarity is learned through social pruning. We model this as a three-stage process:
\begin{enumerate}
    \item \textbf{Innate (Pre-training):} High-entropy, dense representations.
    \item \textbf{Adaptive (RLHF):} Reward-based pruning of the latent space.
    \item \textbf{Specialized (MoE):} Architectural modularization for energy efficiency.
\end{enumerate}

\section{2. Theoretical Framework}

\subsection{2.1 Mathematical Foundations of Hybrid Layers}
Our architecture, \textit{Cortex-13}, integrates the global attention of Transformers with the linear-time history of State Space Models.

\subsubsection{2.1.1 Attention & Hybrid Scaling}
Unlike pure Transformers, our architecture follows the \textit{Jamba v2} specifications. We utilize a 1:7 ratio of Attention to Mamba layers. To maximize context compression, Attention layers explicitly disable Rotary Position Embeddings (RoPE), relying on the implicit sequence tracking of the preceding Mamba blocks.

\subsubsection{2.1.2 Mamba-2 Structured State Space (SSD)}
For the "Stream of Thought" ($h_t$), we upgrade to the \textit{Mamba-2 SSD} algorithm. The state update is formalized through unstructured multiplication in the dual space (SSD):
\begin{equation}
    M = L \circ \Delta \circ (C^\top B)
\end{equation}
Where $M$ represents the attention-equivalent kernel. This formulation allows $O(N)$ linear scaling while maintaining the expressivity of higher-order interactions. The state $h_t$ acts as a thermodynamic bottleneck, forcing the model to \textit{compress} concepts into sparse symbols.

\subsection{2.2 Sparse Autoencoders (SAEs)}
To resolve polysemanticity, we train an SAE on the activation space $x$:
\begin{equation}
    \mathcal{L}_{SAE} = ||x - \hat{x}||^2_2 + \lambda ||f||_1
\end{equation}
Where $\lambda ||f||_1$ enforces sparsity. The resulting directions $f$ represent disentangled "concepts."

\section{3. Methodology}
We implemented \textit{Cortex-13} in PyTorch, featuring:
\begin{itemize}
    \item \textbf{Blocks:} Jamba Hybrid (1:7 ratio). Standard v2 uses 16 layers; for CPU-constrained environments, we utilize a "Micro-Jamba" configuration (3--6 layers).
    \item \textbf{Normalization:} RMSNorm for activation magnitude stability.
    \item \textbf{MoE:} Granular Mixture of Experts (64 Experts, Top-6 routing) with shared experts to prevent knowledge fragmentation.
    \item \textbf{Probing:} Custom \texttt{GlassBoxProbe} for monitoring SSD State Energy and MoE Routing Entropy.
\end{itemize}

The training curriculum features an \textbf{Atomic SFT Warm-up} (3 epochs) to align basic physical syntax, followed by \textbf{RLVR (Reinforcement Learning with Verifiable Rewards)} to establish causal logic.

\section{4. Experiments \& Findings}

\subsection{4.1 RQ1: The Sparsity Transition}
We monitored the $L_0$ norm (number of non-zero elements) of the SAE features throughout training.
\begin{itemize}
    \item \textbf{Stage 1:} Average active features $\approx 450$ per token. High polysemanticity.
    \item \textbf{Stage 2:} Sharp drop to $\approx 35$ active features.
\end{itemize}
\textbf{Finding:} RLHF acts as a "Information Bottleneck," forcing the model to discard noise and settle on discrete concepts.

\subsection{4.2 RQ2: Modular Specialization Dynamics}
We analyzed the Gating Network entropy in the MoE layers.
\begin{equation}
    H(G) = - \sum p_i \log p_i
\end{equation}
Initially $H(G)$ was high (random routing). By Stage 3, $H(G)$ approached 0 for specific syntactic tasks, indicating that Expert \#2 had become the dedicated "Adjective-Noun" processor.

\subsection{4.3 RQ3: Emergence of Symbolic Induction}
We traced the "Induction Heads" (circuit: $A...B \rightarrow A...B$). 
\textbf{Before:} Attention weights were diffuse ($<0.1$).
\textbf{After:} Weights crystallized to $>0.99$. 
This transition transforms probabilistic association into \textbf{Deterministic Logic} ($A \implies B$).

\section{5. Conclusion}
This study confirms that "Glass Box" interpretability is achievable not by simplifying architectures, but by understanding their developmental trajectory. The "Black Box" is simply an untrained brain. Through the lens of Sparse Autoencoders, we have shown that symbolic reasoning is the inevitable low-energy state of a well-aligned neural network.

\section*{Declarations}

\subsection*{Competing Interests}
The authors declare no competing financial or non-financial interests.

\subsection*{Data Availability}
All code, training scripts, and experimental data are publicly available at: \url{https://github.com/nanoglass-core}. The verification suite (\texttt{verify\_all.py}) allows full replication of all empirical claims.

\subsection*{Pre-registration}
This exploratory study was not pre-registered. The experimental protocol evolved iteratively alongside theoretical development, consistent with the discovery-oriented nature of interpretability research.

\subsection*{Author Contributions}
M.R. conceived the theoretical framework, implemented the architecture, and conducted all experiments. A.G. (Antigravity AI) provided supervision and methodological guidance.

\printbibliography

\end{document}
